An observer sees a state of the pendulum near its upper equilibrium through Gaussian noise of width $\sigma$, and the resolution gap $g$ measures how much of the true expansion rate the blurred description loses.
It has been conjectured that $g$ never falls below $c'/\sqrt{\ell+C'}$, $\ell=\ln(1/\sigma)$, whatever the state; so far the lower side was known only for states built from a class of windows.
Here we prove a lower bound for every state, with no condition on its shape.
The only assumptions are a finite first moment and the confinement of the problem itself, $\Var(X_z)\le\pi^2$.
For every such state and every $\sigma\le1.5\times10^{-14}$,
\[
g\;\ge\;\frac{3}{4}\,\frac{\alpha}{\alpha\sigma+1268\,\kb},\qquad \kb^2=289\,\frac{\Lb^{5/2}}{(\ln2\Lb)^{3/2}}+4\Lb+3.57,
\]
with $\alpha=1/(2(\pi^2+\sigma^2))$ and $\Lb=2\ell+\ln(1/\alpha)$; for $\sigma\le10^{-15}$ this gives $g\ge1.7\times10^{-6}\,(\ln2\Lb)^{3/4}/\Lb^{5/4}$ \proofst.
The present power $\ell^{-5/4}(\ln\ell)^{3/4}$ is below the conjectured $\ell^{-1/2}$, the constants are small and the range of $\sigma$ is narrow; we state all three plainly.
Three steps carry the proof.
A three-circle argument in complex time bounds the second-order spread of the posterior for every state: $\keff^2\le289\,L^{5/2}/(\ln2L)^{3/2}+4L+3.57$ with $L=\ln(1/A)$, $A=\sigma^2\Jww$.
Outside an open set $\Oset$ where the $w$-information is of order $\sigma^2$, the remainder of the gap formula is at most half of its main terms.
Inside $\Oset$, an expansion of the gap around the observed point instead of the posterior mean has a remainder of order $\sigma^2\sqrt{\Jww+\Jvv}$, and the same complex-time argument in the isotropic flow controls the one second-order quantity left.
Every constant was recomputed with outward-rounded interval arithmetic, and every joint of the chain was evaluated on $6062$ recorded states without a failure \numst.
We state what is proved, derived and numerical, and what remains open: the power $\ell^{-1/2}$ and the range $\sigma>1.5\times10^{-14}$, with the inequality each rests on.
